\section{The nuclear quartet: two commuting nuclei}

This section records claim nodes and dependency edges.

\begin{proposition}[Heyting algebra]\label{thm:heyting}\lean{Obsmetrologia.HeytingAlgebra}\leanok\uses{thm:floor}
\texttt{HeytingAlgebra} is a meet-semilattice (commutative, associative, idempotent \texttt{meet}) with a relative pseudocomplement \texttt{imp} adjoint to meet via the Heyting adjunction $b \sqcap a \leq c \leftrightarrow b \leq (a \Rightarrow c)$, the order being $a \leq b := a \sqcap b = a$. Its induced order is reflexive, transitive and antisymmetric, meet is monotone, and modus ponens holds --- the algebraic floor on which nuclei act.
\end{proposition}
\begin{proposition}[Nucleus]\label{thm:nucleus}\lean{Obsmetrologia.Nucleus}\leanok\uses{thm:heyting}
A \texttt{Nucleus} on a Heyting algebra is an inflationary ($a \leq j\,a$), idempotent ($j\,(j\,a) = j\,a$), meet-preserving ($j\,(a \sqcap b) = j\,a \sqcap j\,b$) endomap. Monotonicity follows from meet-preservation. A nucleus is a closure operator that respects the Heyting structure, not merely the lattice.
\end{proposition}
\begin{theorem}[Fixed points are closed under $\Rightarrow$]\label{thm:fix-imp}\lean{Obsmetrologia.Nucleus.fix_imp}\leanok\uses{thm:nucleus}
\texttt{Nucleus.fix\_imp} proves that if $b$ is $j$-fixed, so is $a \Rightarrow b$. This is why meet-preservation matters: it makes $\mathrm{Fix}(j)$ a Heyting \emph{sub-algebra}, not merely a sub-lattice. Without the meet law the fixed points would reflect only the lattice order.
\end{theorem}
\begin{theorem}[Reflected joins absorb reflected input]\label{thm:nucleus-absorb-join}\lean{Obsmetrologia.Order.Nucleus.absorb_join}\leanok\uses{thm:nucleus}
For a nucleus $j$, the reflected join $j(j\,a \vee b)$ equals $j(a \vee b)$.
Once the left input has already been sent through the reflector, joining and
reflecting again sees no extra information. This is the active dependency-free
form of the archive's join-absorption nucleus law.
\end{theorem}
\begin{theorem}[Reflected meet-joins absorb reflected inputs]\label{thm:nucleus-absorb-join-meet}\lean{Obsmetrologia.Order.Nucleus.absorb_join_meet}\leanok\uses{thm:nucleus-absorb-join}
For a nucleus $j$, the reflected join
$j((j\,a \wedge j\,b)\vee c)$ equals $j((a\wedge b)\vee c)$.
Meet preservation first rewrites $j\,a\wedge j\,b$ as $j(a\wedge b)$;
the join-absorption law then removes the redundant reflection. This is the
active critical-pair closure carried by the old nuclear reduction notes.
\end{theorem}
\begin{theorem}[Dynamic residuation]\label{thm:nucleus-dynamic-residuation}\lean{Obsmetrologia.Order.Nucleus.galois_connection_meet}\leanok\uses{thm:nucleus, thm:fix-imp}
For a nucleus $j$, reflection after meet is adjoint to implication out of the
reflected left input:
$j(a\wedge b)\leq j(c)$ iff $j(b)\leq (j(a)\Rightarrow j(c))$.
The result uses only meet preservation and the Heyting adjunction; no ambient
completeness or maximum construction is assumed.
\end{theorem}
\begin{theorem}[Dynamic residuation is stable]\label{thm:nucleus-dynamic-residuation-stability}\lean{Obsmetrologia.Order.Nucleus.galois_connection_stability}\leanok\uses{thm:nucleus-dynamic-residuation}
If $j(a\wedge b)\leq j(c)$ and $j(c\wedge d)\leq j(e)$, then
$j((a\wedge b)\wedge d)\leq j(e)$. This is the proved compositional stability
of the dynamic-residuation relation in the active order hierarchy.
\end{theorem}
\begin{theorem}[Stable envelope lands in the joint fixed corner]\label{thm:nucleus-stable-envelope-fixed}\lean{Obsmetrologia.Order.CommutingNuclei.stableEnvelope_fix_pi}\leanok\uses{thm:nuclear-square, thm:fix-imp}
For two commuting nuclei that both fix bottom, the stable-envelope candidate
$Env(a)=i_{j_1}(a)\wedge i_{j_2}(a)$, with $i_j(a)=j(a\Rightarrow \bot)\Rightarrow\bot$,
is fixed by the composite nucleus. The proof uses the fixed-point closure of
Heyting negation under a bottom-preserving nucleus. The stronger claim that
$Env(a)$ is the greatest joint fixed point below $a$ is not asserted here.
\end{theorem}
\begin{proposition}[Nuclear formula syntax]\label{thm:nuclear-formula-syntax}\lean{Obsmetrologia.Positing.NuclearFormula}\leanok\uses{thm:nucleus}
\texttt{NuclearFormula} is the dependency-free Heyting-modal propositional
syntax with variables, $\top$, $\bot$, conjunction, disjunction, implication,
and the two modal constructors \texttt{Saturate} and \texttt{Transfer}. This is
the active syntax carrier extracted from the old fiber nuclear language.
\end{proposition}
\begin{proposition}[Nuclear formula interpretation]\label{thm:nuclear-formula-interpretation}\lean{Obsmetrologia.Positing.NuclearFormula.interp}\leanok\uses{thm:nuclear-formula-syntax, thm:nuclear-square}
\texttt{NuclearFormula.interp} interprets the syntax in any Heyting algebra
equipped with two commuting nuclei: the Heyting connectives go to the ambient
operations, \texttt{Saturate} goes to $j_1$, and \texttt{Transfer} goes to $j_2$.
\end{proposition}
\begin{proposition}[Nuclear formula step grammar]\label{thm:nuclear-formula-step}\lean{Obsmetrologia.Positing.NuclearFormula.Step}\leanok\uses{thm:nuclear-formula-interpretation, thm:nucleus-absorb-join-meet}
\texttt{NuclearFormula.Step} records the archive's oriented root rules R1--R23
and congruence closure for rewriting under modal and Heyting constructors. It
is a grammar of the proposed normalization system, not by itself a confluence
or canonical-form theorem.
\end{proposition}
\begin{theorem}[Nuclear rewrite measures]\label{thm:nuclear-formula-measures}\lean{Obsmetrologia.Positing.NuclearFormula.entanglement_R7}\leanok\uses{thm:nuclear-formula-step}
\texttt{entanglement\_R7} proves that the commutation rule
\texttt{Saturate (Transfer M) -> Transfer (Saturate M)} strictly decreases the
Saturate-inside-Transfer entanglement measure. The companion green lemmas
\texttt{nuclearW\_R3\_decreases}, \texttt{nuclearW\_R4\_decreases}, and
\texttt{nuclearW\_R7\_preserved} establish the polynomial-weight side of the
archive's termination measure.
\end{theorem}
\begin{theorem}[Nuclear canonical forms frontier]\label{thm:nuclear-formula-canonical-frontier}\notready\uses{thm:nuclear-formula-measures}
The archive sketches full strong normalization, local confluence, and canonical
forms for the R1--R23 system, but its canonical-form file leaves many local
confluence cases as \texttt{sorry}. The active frontier is to prove those
critical pairs dependency-free, then derive confluence and uniqueness of normal
forms.
\end{theorem}
\begin{proposition}[Nuclear square]\label{thm:nuclear-square}\lean{Obsmetrologia.NuclearSquare}\leanok\uses{thm:nucleus}
A \texttt{NuclearSquare} on a Heyting algebra is a pair of nuclei $j_1, j_2$ together with the commutation law $j_1\,(j_2\,a) = j_2\,(j_1\,a)$. The pair is read either as the two orders $(j_k, j_t)$ of the bilattice or as the temporal pair $(\sigma, \Delta)$ --- saturation and transfer.
\end{proposition}
\begin{theorem}[Nuclear Quartet Theorem]\label{thm:nuclear-quartet}\lean{Obsmetrologia.NuclearSquare.compNucleus}\lean{Obsmetrologia.Order.CommutingNuclei.transferring_without_saturating_not_pi_fixed}\lean{Obsmetrologia.Order.CommutingNuclei.transferring_without_saturating_pi_ascends}\lean{Obsmetrologia.Order.NuclearQuartet.active_not_stable_fixed_of_not_potential}\lean{Obsmetrologia.Order.NuclearQuartet.active_not_potential_stable_ascends}\leanok\uses{thm:nuclear-square, thm:fix-imp}
The joint closure $\pi = j_1 \circ j_2$ is itself a nucleus: \texttt{NuclearSquare.compNucleus} constructs it, with idempotency turning on the commutation law. Its companion \texttt{NuclearSquare.fix\_pi\_iff} proves $\mathrm{Fix}(\pi) = \mathrm{Fix}(j_1) \cap \mathrm{Fix}(j_2)$ --- the joint-quiescent corner is exactly the joint fixed point, the hinge again being commutativity. \texttt{NuclearSquare.square} restates the commuting square: both paths $A \to A^{*}$ agree. In the standalone \texttt{Order.CommutingNuclei} register, Lean also names the contrapositive boundary: a $j_2$-fixed point outside $\mathrm{Fix}(j_1)$ is not fixed by $\pi$, and extensivity of the composite expresses that failure as ascent $a\le\pi(a)$ with $\pi(a)\ne a$. The quartet restatement says the same thing as: stable but not active is not absurd-fixed, and the absurd/joint corner raises it.
\end{theorem}
\begin{proposition}[The four corners of the quartet]\label{thm:quartet-corner}\lean{Obsmetrologia.NuclearSquare.corner}\leanok\uses{thm:nuclear-quartet, thm:four-object}
\texttt{NuclearSquare.corner} reads the four sub-algebras off the four articulation phases: absurd $\mapsto$ the free carrier (neither nucleus), potential $\mapsto j_1$-quiescent, active $\mapsto j_2$-quiescent, stable $\mapsto$ joint-quiescent ($\pi$). The quartet's corners are the recurring FOUR, lifted from values to fixed-point sub-algebras.
\end{proposition}

This section records claim nodes and dependency edges.

\begin{proposition}[Family of nuclear squares]\label{thm:square-family}\lean{Obsmetrologia.SquareFamily}\leanok\uses{thm:nuclear-square}
A \texttt{SquareFamily} over an index type $I$ assigns a \texttt{NuclearSquare} to each fiber (history or step). The nuclear pair $(j_1^i, j_2^i)$ may differ per fiber --- this is the recurrence: not one square but a square at every fiber.
\end{proposition}
\begin{proposition}[The family varies]\label{thm:varies}\lean{Obsmetrologia.SquareFamily.Varies}\leanok\uses{thm:square-family}
\texttt{SquareFamily.Varies} asserts that some two fibers carry genuinely different nuclear pairs. Constant families leave automorphisms invisible; a varying family is what makes history-separation (aperiodicity) provable. \texttt{SquareFamily.fiber\_quartet} confirms that every fiber realizes the Nuclear Quartet Theorem regardless of variation.
\end{proposition}
\begin{example}[A concrete varying family]\label{ex:varying-family}\lean{Obsmetrologia.Charge.varyingFamily}\leanok\uses{thm:varies, thm:four-bilattice}
On the charge knowledge-Heyting-algebra \texttt{Charge.heytingK} (with \texttt{Charge.joinNucleusK} the nucleus raising toward a bound by $\vee_k$), \texttt{Charge.varyingFamily} is a two-fiber family: fiber \texttt{true} raises toward $t$, fiber \texttt{false} toward $\top_k$. \texttt{Charge.varyingFamily\_varies} proves the $j_1$ nuclei genuinely differ --- a concrete bigraded-diamond recurrence over the finite charge bilattice.
\end{example}
